\section{Introduction}
\label{sec.intro}

Bitcoin tolerates an adversary holding just under half of the mining power.
That threshold is not free. It is bought with an assumption about the network:
a message an honest validator sends at round $i$ reaches every other honest
validator at round $i+1$. Weaken the assumption and the threshold falls. Under
a bounded delay the tolerable adversary shrinks as the delay
grows~\cite{pss16,dembo20}, and under partial synchrony no protocol reaches one
half in both safety and liveness at once~\cite{cuamdr}. Synchrony among the
validators is load bearing, and the backbone analysis of Bitcoin assumes it
throughout~\cite{gkl24}.

Validators are not the only parties who care about the answer. A client does
not mine and does not run the protocol. It wants one thing, which is to read
the ledger and to be sure that what it reads will not later be undone. Wallets
are clients. So is a payment processor checking that a transfer went through,
and so is anyone who looks at the chain without maintaining it. These parties
are not on the validators' network and have no reason to be. A phone on an
aeroplane is a client. So is a laptop behind a router someone else controls.
Many such parties sidestep the problem by asking a hosted node to read the
chain on their behalf, which does not answer the question so much as move it to
whoever runs the service.

The decentralised alternatives face the question squarely, and every one of
them answers it the same way.
Section 8 of~\cite{nakamoto08} tells a client to keep the headers of the longest
proof of work chain, ``which he can get by querying network nodes until he's
convinced he has the longest chain''. The quoted clause is an assumption written
as advice. It asks that one of the nodes the client reaches be honest, and it
asks that the chain that node holds be current. Superblock proofs~\cite{nipopow},
probabilistic sampling~\cite{flyclient} and the interactive protocol
of~\cite{blink} shrink what the client downloads, in the last case to a constant,
and each of them inherits the clause untouched. All of them need one honest party
among those the client speaks to, and all of them need the network to deliver on
time. A client that has been eclipsed~\cite{eclipse} has neither.

This paper asks what a client can do with neither, and the answer is that it
can still confirm a prefix and can work out which one. The observation is about
chain growth. Honest validators extend the chain at a rate fixed by the mining
target and by how much of the mining power is honest, and a client knows both
of these before it connects to anything. Knowing the round, it therefore knows
a length the honest chain has already passed. A chain shorter than that was not
produced by honest validators keeping pace. Its recent blocks were withheld
while it was built, and the rate at which the adversary mines bounds how many
such blocks there can be. A shortfall in length is not a sign that the client
is merely behind and holding something older and safer. It is a measurement of
how much work could have gone into a private fork.

Two bounds meet, and the rule falls out of them. Write $f$ for the probability
that some honest validator mines in a round, $r$ for the round at which the
chain reaches the client and $\hgt$ for its height. No chain has height above
$\frac{n}{n-t}fr$ at round $r$, so a chain that is short forked from the honest
chain early. The adversary mines about $\frac{t}{n-t}f$ blocks per round, so a
fork that began early had that many rounds in which to grow. The two pull
against each other, and eliminating the round of the fork between them leaves
\[
    \prn(r) = \frac{t}{n-t}\left(\frac{n}{n-t}fr-\hgt\right)
\]
blocks, which is what the client discards. \Cref{fig.confirmed} is the same
statement drawn in the plane, and \Cref{def.client} is the rule with the
corrections a typical execution requires. Nothing in it refers to a peer. The
client evaluates it on whatever chain arrives, from whoever sends it, and the
prefix that survives is a prefix of the chain every honest validator holds.

Three things about the rule deserve stating plainly, because each runs against
the way confirmation is usually described. The first is that a shorter chain is
worse than a longer one, not safer. The reading that a client behind the tip is
merely looking at older and therefore settled history is wrong, since honest
validators are never behind, and every block of shortfall buys the adversary
$\frac{t}{n-2t}$ blocks of fork. That exchange rate crosses one at $t=n/3$.
The second is that no constant confirmation depth is sound. A chain that an
honest validator adopted at round $r$ and handed over in the same round still
has $\frac{t^2}{(n-t)^2}fr$ blocks pruned from it, and that quantity grows
without bound as the execution runs. Six confirmations, or six hundred, fail
once the chain is old enough. The third is that liveness survives this anyway.
The depth grows, but the fraction of the chain it removes does not: it stays at
$\frac{t^2}{(n-t)^2}$, the square of the adversary's share of the queries
relative to the honest share. Against an adversary with a tenth of the mining
power that fraction is about $1.2\%$, so a client handed a chain of one million
blocks discards the last twelve thousand of them. Against a quarter it is
$11\%$ and against a third it is $25\%$. The confirmed prefix grows at a
constant rate, and a client that runs the rule again later confirms strictly
more.

Our client is not a sleepy client, and the difference is the point. A sleepy
party in the sense of~\cite{sleepy} knows when it was asleep, receives on waking
everything it missed, and assumes synchrony while it is awake. The client models
of~\cite{cuamdr} inherit that shape, and what sleeping costs there is the
ability to time the arrival of a message. Our client keeps its clock and gives up
everything else. It assumes no bound on the delay of its own link, at any round,
and it assumes that no party it speaks to is honest. The clock is all that is
left, and counting rounds since the execution began turns out to be enough.

Clients that never speak to one another still agree. The output of one is a
prefix of the output of the other, whatever each was sent and whenever each
joined, even if one of them is eclipsed and the other is not. This is the
guarantee a reader of a replicated state machine expects, and here it costs no
assumption about either client's network. Finally, the bound cannot be
improved. Against an adversary that mines in private, a client that confirmed
any more would confirm a block that the honest validators go on to abandon, so
the prefix our rule returns is the longest prefix any client can safely return.

\paragraph{Contributions}

We work in the static synchronous backbone model of~\cite{gkl24} and change
what a client is allowed to assume: nothing about its own network, and nothing
about the honesty of the parties that answer it.
\begin{enumerate}
    \item We give a confirmation rule that a client evaluates from the round
        and the chain in front of it (\Cref{sec.construction}).
    \item We prove that what it returns is a prefix of every chain an honest
        validator adopts from that round onwards, and that any two clients
        running it return prefixes of one another (\Cref{sec.safety}).
    \item We prove it live, in one bound that covers an honest counterparty and
        an adversarial one, and we show the confirmed prefix grows at a constant
        rate (\Cref{sec.liveness}).
    \item We prove the bound tight by exhibiting a private mining adversary
        that attains it, and we show that no constant confirmation depth is
        sound (\Cref{sec.tightness}).
\end{enumerate}
\Cref{sec.related} places the work among existing client protocols.
\Cref{sec.model} fixes the model, \Cref{sec.construction} gives the
construction, \Cref{sec.analysis} proves it safe and live,
\Cref{sec.tightness} proves it tight, and \Cref{sec.discussion} collects the
remarks.
